\paragraph{\trEquivariant{}\trTransversality。}\label{sub:equivariant-transversality}
我们的下一个课题，是把 $\ZZ/2$-对称实现为链层面上的一个朴素群作用。我们不以最大
普遍性为目标，而是选择施加一列几何约束，以简化解法。
\begin{assumption}\label{thm:fixed-point-set-assumption}
在我们给定的带角\trExactSymplectic\trManifold{} $M$ 与\trInvolution{} $\iota_M$ 中，
我们要求：\trFixedPoint{}集 $M^{\mathrm{fix}}$ 的实\trCodimension{}为二，并且其\trNormalBundle{}
$\nu$ 是平凡的。我们实际上想取定该丛的一个复\trTrivialization{} $\nu\to\CC$，并附加
如下的\textbf{\trPartialIntegrability}（partial integrability）条件：该\trTrivialization{}应当是某个映射 $\pi\colon U\to\CC$ 的%
\trDifferential，其中 $U\subset M$ 是\trFixedPoint{}集的某个邻域；$\pi$ 满足 $\pi\circ\iota_M=-\pi$，
并且 $\bar{\partial}\pi$（相对于 $I_M$ 取）沿 $M^{\mathrm{fix}}$ 一阶消失。
\end{assumption}
设我们给定 $I_M$ 与\trTrivialization{} $D\pi|M^{\mathrm{fix}}$。若 $N(I_M)$ 是 Nijenhuis 张量，则
具备上述性质的任何 $\pi$ 之存在性蕴含 $D\pi\circ N(I_M)$ 在 $M^{\mathrm{fix}}$ 上为零，
这解释了“\trPartialIntegrability”这个名字。
\begin{assumption}\label{thm:normal-angle-assumption}
在一切 $\iota_M$-不变的\trExact\trLagrange\trSubmanifold{} $L$ 之中，我们只考虑这样
一些：它们或者根本不同\trFixedPoint{}集相交，或者从一个常值方向趋近它。这意味着存在某个
$\sigma\in\RR$——$L$ 的\textbf{\trNormalAngle}（normal angle）——使得
\begin{equation*}
D\pi(TL_x)=e^{i\sigma}\RR\subset\CC
\end{equation*}
对一切 $x\in L^{\mathrm{fix}}$ 成立。$\sigma$ 的可能选取有一个 $\ZZ$ 的余地，但我们将
假定：对每个相关的\trLagrange\trSubmanifold，一个优先的选取已经做出。
\end{assumption}
在进入下一部分之前，我们需要引进一些术语。设 $H$ 是 $M$ 上一个 $\iota_M$-不变的
函数。在每一点 $x\in M^{\mathrm{fix}}$ 处，$D^2H$ 的法向部分是 $\nu_x$ 上一个\trWellDefined{}的%
\trQuadraticForm（因为一阶导数必定为零）。若对 $Z\in\nu_x$ 有
\begin{equation}\label{eq:normal-hessian-rotation}
(D^2H)(Z)=h|Z|^2
\end{equation}
——其中 $|\cdot|$ 是把 $\nu$ 等同于 $TM|M^{\mathrm{fix}}$ 的反不变部分所得的%
\trHermitian{}\trMetric{}——则称 $H$ \textbf{无穷小旋转}（infinitesimally rotate）$\nu_x$，
其速度为 $h\in\RR$。\trEquivalence{}地，取\trHamiltonianVectorField{} $X_H$，用 $D\pi$
去\trProjection{}它，得到一个函数 $D\pi(X_H)\colon U\to\CC$。于是 \eqref{eq:normal-hessian-rotation}
说的是：在 $x$ 附近的一阶范围内，$D\pi(X_H)=ih\pi$。
\begin{assumption}\label{thm:floer-datum-assumption}
取 $(L_0,L_1)$ 为一对如\cref{thm:normal-angle-assumption} 的不变\trExact\trLagrange%
\trSubmanifold，二者都与\trFixedPoint{}集相交。设 $\sigma_0,\sigma_1$ 是它们的\trNormalAngle。
在考虑 $\ZZ/2$-不变的\trFloerDatum{} $(H,J)$ 时，我们想施加下列附加要求。其一，正如
对 $I_M$ 的假定那样，我们想让相对于任何 $J_t$ 取的导数 $\bar{\partial}\pi$ 沿
$M^{\mathrm{fix}}$ 一阶消失。其二，设 $\phi$ 是 $H$ 的“流”。那么，应当存在 $h\in\RR$
满足
\begin{equation}\label{eq:normal-angle-chamber}
\sigma_1-\sigma_0-h\in(-\pi;0)
\end{equation}
并且存在一个开子集 $V\subset[0;1]\times M^{\mathrm{fix}}$，它对一切 $t\in[0;1]$ 都
包含 $\{t\}\times\phi^t(L^{\mathrm{fix}}_0)$，使得 $H(t,\cdot)$ 对一切 $(t,x)\in V$
都以速度 $h$ 无穷小旋转 $\nu_x$。
\end{assumption}
我们退而只在一个足够大的开子集上控制法向 Hessian，这是因为按\trFloerDatum{}的定义，
$H(t,\cdot)$ 必须在 $\partial M$ 附近为零，从而不可能在整个 $M^{\mathrm{fix}}$ 上满足
\eqref{eq:normal-hessian-rotation}。要紧的是：对 $x\in L^{\mathrm{fix}}_0$，
$D\phi^t_x\colon\nu_x\to\nu_{\phi^t(x)}$ 是旋转 $e^{iht}$。特别是，$\phi^1(L_0)$ 从方向
$e^{i(\sigma_0+h)}\RR$ 趋近\trFixedPoint{}集。仍按\trFloerDatum{}的定义，它必须与 $L_1$
\trTransverse，这本身就已蕴含 $\sigma_1-\sigma_0-h\in\RR\setminus\pi\ZZ$（当然，除非
$\phi^1(L_0)\cap L_1\cap M^{\mathrm{fix}}=\emptyset$）。附加要求 \eqref{eq:normal-angle-chamber}
则挑出 $\RR\setminus\pi\ZZ$ 的一个“房”即\trConnectedComponent。这看起来可能相当空洞，
因为总可以改变 $\sigma_0$ 或 $\sigma_1$ 来达到想要的结果；但一旦同时考虑多于一个%
\trFloerDatum，它就有了意义。
具备这些性质的函数仍然比比皆是。特别是，我们既可以任意指定 $H$ 在\trFixedPoint{}集上
的值，也可以任意指定它在该集合的任意小邻域之外的值（当然要受在边界附近为零、
以及 $\iota_M$-不变的限制）。对于\trAlmostComplex{}结构，同样的话也成立，论证如下。首先，把
$M^{\mathrm{fix}}$ 的邻域 $U$ 取小一些（若有必要），使 $\pi|U$ 成为具辛\trFiber{}的\trSubmersion{}。
\trTangentBundle{} $TM|U$ \trSplitting{}为 $TM^v=\ker(D\pi)$ 及其辛\trOrthogonal{}补 $TM^h$。后者被 $D\pi$ \trTrivialization{}，而它
限制到 $M^{\mathrm{fix}}$ 上恰好是 $\nu$。由于 $I_M$ 与\trTrivialization{} $D\pi|M^{\mathrm{fix}}$
\trCompatible，我们知道：沿\trFixedPoint{}集，$I_M$ 的 $(hv)$ 分量为零，而 $(hh)$ 分量等于
$i$。对 $\bar{\partial}\pi$ 的附加要求是说，这两个条件应当扩张到一阶导数（实际上，
就 $(hh)$ 分量而言，这已可由 $\iota_M$-对称性推出）。在寻找具有同一性质的其它%
\trAlmostComplex{}结构时，显然可以随意选取沿 $M^{\mathrm{fix}}$ 的 $(vv)$ 分量。
\begin{lemma}\label{thm:regularity-fixed-set}
设 $L_0,L_1$ 是不变的\trExact\trLagrange\trSubmanifold，$(H,J)$ 是如%
\cref{thm:floer-datum-assumption} 的一个不变\trFloerDatum。取某个 $u\in\scr{M}_Z(y_0,y_1)^*$，
使其像完全含于 $M^{\mathrm{fix}}$ 之内。那么，$u$ 在 $M$ 中\trRegular{}当且仅当它在
$M^{\mathrm{fix}}$ 中\trRegular{}。
\end{lemma}
\begin{proof}
\trLinearization{}算子是 $\ZZ/2$-\trEquivariant{}的，因而\trSplitting{}为
\begin{equation*}
D_{Z,u}=D^-_{Z,u}\oplus D^+_{Z,u}.
\end{equation*}
不变部分易于辨认。具体地说，令 $u^{\mathrm{fix}}$ 表示把 $u$ 看作 $M^{\mathrm{fix}}$ 内%
\trFloerEquation{}的解——相对于\trLagrange\trSubmanifold{} $(L^{\mathrm{fix}}_0,
L^{\mathrm{fix}}_1)$ 与\trFloerDatum{} $(H|M^{\mathrm{fix}},J|M^{\mathrm{fix}})$——那么
\begin{equation}\label{eq:invariant-part-operator}
D_{Z,u^{\mathrm{fix}}}=D^+_{Z,u}.
\end{equation}
至于反不变部分，可以用 $D\pi$ 把相关的\trVectorBundle{}等同于 $E=Z\times\CC\to\CC$，
于是\trSubbundle{} $F\subset E|\partial Z$ 就变成 $F_{s,0}=e^{i\sigma_0}\RR$、
$F_{s,1}=e^{i\sigma_1}\RR$。沿 $M^{\mathrm{fix}}$ 的一阶范围内，我们知道
$D\pi\circ J_t=i\circ D\pi$。此外，当 $|s|\gg 0$ 时 $x=u(s,t)$ 靠近 $\phi^t(L_0)$，
故 $(t,x)\in V$，这意味着 $H(t,\cdot)$ 的\trHamiltonianVectorField{} $X(t,\cdot)$
满足 $D\pi(X(t,x))=ih\pi(x)$（在 $x$ 处取到一阶）。这意味着：在诸\trStriplikeEnd{}上，
\begin{equation}\label{eq:anti-invariant-operator}
D^-_{Z,u}=\partial_s+i\partial_t+h.
\end{equation}
可以把算子\trDeformation{}到使 \eqref{eq:anti-invariant-operator} 处处成立，这就表明
$\mathrm{index}(D^-_{Z,u})=0$，例如参见 \cref{thm:path-space-det-line}。另一方面，
\cref{thm:winding-injectivity}（$\mu=0$、$|\Sigma^-|=1$）说 $D^-_{Z,u}$ 是\trInjective{}
的，因而必\trInvertible{}。鉴于 \eqref{eq:invariant-part-operator}，这就蕴涵了想要的结果。
\end{proof}
有了这个工具，要证明 \cref{thm:floer-datum-assumption} 所定义的那个相当特殊的类
之中的\trFloerDatum{}是\trGeneric{}\trRegular{}的，就成了例行公事。给定一个这样的数据，先只
看\trFixedPoint{}集 $M^{\mathrm{fix}}$，并对它稍加\trPerturbation{}，使之在那里\trRegular{}。由
\cref{thm:regularity-fixed-set}，这保证了 $M$ 中一切不变\trFloer{}\trTrajectory{}的\trRegularity{}。其余
的\trTrajectory{}，可以由标准\trTransversality{}论证的一个版本做成\trRegular{}的；参见 [85, Proposition 5.13]。
设 $(L_0,\ldots,L_d)$（$d\ge 2$）是如\cref{thm:normal-angle-assumption} 的 $\iota_M$-不变%
\trExact\trLagrange\trSubmanifold，每一个都与\trFixedPoint{}集相交。设 $\sigma_0,\ldots,\sigma_d$
是它们的\trNormalAngle。设我们已为 $(L_0,L_d)$ 与 $(L_{k-1},L_k)$（$1\le k\le d$）
诸对选好了如\cref{thm:floer-datum-assumption} 的\trFloerDatum。考虑 $(d{+}1)$ \trPointedDisc{}
的\trUniversal\trFamily{} $\scr{S}^{d+1}\to\scr{R}^{d+1}$，配以\trLagrange\trLabel{}
$(L_0,\ldots,L_d)$，并给它配备某个 $\iota_M$-不变的\trPerturbationDatum{} $(K,J)$，与
前面\trFloerDatum{}的选取\trCompatible。
\begin{lemma}\label{thm:regularity-polygons}
设 $(r,u)\in\scr{M}^{d+1}(y_0,\ldots,y_d)=\scr{M}_{\scr{S}^{d+1}}(y_0,\ldots,y_d)$ 是
一个\trInhomogeneous{}\trPseudoholomorphic\trPolygon，其像完全含于 $M^{\mathrm{fix}}$ 之内。那么，$u$ 在 $M$ 中%
\trRegular{}当且仅当它在 $M^{\mathrm{fix}}$ 中\trRegular{}。
\end{lemma}
\begin{proof}
与前面一样，我们要证算子 $D_{\scr{S}^{d+1},r,u}$ 的反不变部分\trInvertible{}。第一步同样是把
该部分写成 $D^-_{S,u}\colon W^{1,p}(S,E,F)\to L^p(S,E)$，其中 $S=\scr{S}^{d+1}_r$、
$E=S\times\CC$，而 $F|C_k=e^{i\sigma_k}\RR$。在每个\trStriplikeEnd{}上，该算子的性态
由\trFloerDatum{}规定，也就是说 $\epsilon^*_{S,k}D^-_{S,u}=\partial_s+i\partial_t+h_k$，
如 \eqref{eq:anti-invariant-operator}，其中诸 $h_k$ 满足合适的条件
\eqref{eq:normal-angle-chamber}。然后组合 \cref{thm:general-surface-index-formula}
（它给出\trIndex{}的消失性）与 \cref{thm:winding-injectivity}。
\end{proof}
与 \cref{thm:regularity-fixed-set} 之后的讨论平行，这蕴含：$\iota_M$-不变%
\trPerturbationDatum{} $(K,J)$ 的一个\trGeneric{}选取，使诸\trModuliSpace{}
$\scr{M}_{\scr{S}^{d+1}}$ 全都\trRegular{}（事实上，这一回的论证还要简单些，因为 $(d{+}1)$
\trPointedDisc{}没有\trAutomorphism）。
\begin{remark}\label{thm:weaker-assumptions-remark}
如 \cref{thm:zeros-finite-multiplicity} 中所说，可以证明 \cref{thm:winding-injectivity}
的一个版本，它对相关算子在诸\trStriplikeEnd{}上的性态不需要这么强的控制。那反过来
使人能够把这里施加的某些条件换成纯粹\trTopology{}的对应物——只要它们强到足以保证反不变%
\trLinearization{}算子\trIndex{}为零即可。换个方向，也可以考虑\trFixedPoint{}集\trCodimension{}更高的 $\ZZ/2$-作用，
但那时对反不变算子的显式控制必须保持，而且事实上要扩张到整条\trRiemann\trSurface{}
上，因为维数论证不再可用。
\end{remark}
